% Propietario conservado: /Users/ruben/Documents/ChatGPT/jueces y controles/REVISION_ARGUMENTAL_APERTURAS_CIERRES_20260915/IV/ES/sections/iv_apendice_enumeraciones.tex
% RESULTADO_RECUPERADO. Fuentes conservadas en la entrega:
% antecedentes/controles_articulo_I/FIRMAS_EMISION.json:
% plus_classes, times_classes, seed_id_formula y coordinate_convention.
% antecedentes/controles_articulo_I/procedencia_excepcional.md:
% sección «Comprobación finita ejecutada», matriz A y seis generadores.
% Se explicitan las enumeraciones ya utilizadas; no se cambia su dominio.
\section{Enumeraciones finitas de la emisión y de las estrellas de Witt}
\label{iv:iv:apendice-enumeraciones}

Las dos enumeraciones de este apéndice corresponden a objetos distintos.
La primera clasifica las firmas de los cursores de la emisión APP;
la segunda calcula el grupo generado por las estrellas del diseño de Witt.
Se conservan sus dominios y reglas, de modo que cada tabla pueda
reproducirse sin utilizarla como entrada del cálculo.

\subsection{Índices de las condiciones iniciales y recurrencia de las firmas}
\label{iv:iv:apendice-firmas}

Sean \(i,j\in\{0,\ldots,8\}\) y \(a\in\{0,1,2,3\}\), con el orden
de direcciones \((N,E,S,O)\) y los vectores
\[
 v_N=(-1,0),\quad v_E=(0,1),\quad
 v_S=(1,0),\quad v_O=(0,-1).
\]
El índice
\[
 \operatorname{id}(i,j,a)=4(9i+j)+a
\]
recorre exactamente \(\{0,\ldots,323\}\). Su inversa se obtiene
mediante \(a=\operatorname{id}\bmod4\),
\(j=\lfloor\operatorname{id}/4\rfloor\bmod9\) e
\(i=\lfloor\operatorname{id}/36\rfloor\).
Dos índices, uno por hoja, enumeran las \(324^2\) condiciones conjuntas
de \eqref{iv:eq:semillas}; no se identifica ese dominio con las 81 posiciones.

Para cada índice se inicia el cursor en \((i,j)\), con vector \(v_a\).
Se recorren los instantes \(t=1,\ldots,54\), repartidos en seis ventanas
de nueve instantes. Escribamos \(r_t=1+((t-1)\bmod9)\). La regla es:
\begin{enumerate}
\item En la hoja aditiva se actúa cuando \(r_t\in\{1,4,7\}\):
se añade \(\rho_9(i+j+2)\) al acumulador de la ventana y se
desplaza el cursor por \(v_a\), módulo nueve.
\item En la hoja multiplicativa se actúa cuando \(r_t\in\{2,5,8\}\):
se añade \(\ell_9(\rho_9((i+1)(j+1)))\) y se desplaza el cursor
por \(v_a\), módulo nueve. Se conserva la extensión por cero de
\(\ell_9\) fuera de las unidades, declarada en la sección~\ref{iv:sec:extension}.
\item En los demás instantes ese cursor no se desplaza. Después
de los instantes \(27\) y \(54\) se reemplaza \(v_a\) por \(-v_a\).
\item Al cerrar cada ventana se registra su acumulador módulo diez
y se reinicia únicamente el acumulador; no la posición ni la orientación.
\end{enumerate}
Esto define dos funciones sobre los 324 índices,
\[
 s:\{0,\ldots,323\}\longrightarrow\{0,\ldots,9\}^6,\qquad
 e:\{0,\ldots,323\}\longrightarrow\{0,\ldots,9\}^6.
\]
La multiplicidad de una firma es el número de índices con esa imagen.
El índice testigo de las tablas es el menor de su fibra.
Las preimágenes completas se recuperan aplicando la recurrencia a los
324 índices y conservando los que producen la firma indicada.

\begin{longtable}{rcrr}
\caption{Firmas aditivas: imagen de \(s\), multiplicidad e índice testigo.}
\label{iv:iv:tabla-firmas-aditivas}\\
\toprule
Clase & Firma & Multiplicidad & Índice testigo\\
\midrule
\endfirsthead
\toprule
Clase & Firma & Multiplicidad & Índice testigo\\
\midrule
\endhead
1&122564&18&17\\
2&122898&18&24\\
3&212645&18&5\\
4&212988&18&0\\
5&221456&18&29\\
6&221889&18&12\\
7&456221&18&28\\
8&465898&18&21\\
9&546988&18&9\\
10&564122&18&16\\
11&645212&18&4\\
12&654889&18&33\\
13&889221&18&13\\
14&889654&18&32\\
15&898122&18&25\\
16&898465&18&20\\
17&988212&18&1\\
18&988546&18&8\\
\bottomrule
\end{longtable}

\begin{longtable}{rcrr}
\caption{Firmas multiplicativas: imagen de \(e\), multiplicidad e índice testigo.}
\label{iv:iv:tabla-firmas-multiplicativas}\\
\toprule
Clase & Firma & Multiplicidad & Índice testigo\\
\midrule
\endfirsthead
\toprule
Clase & Firma & Multiplicidad & Índice testigo\\
\midrule
\endhead
1&000000&108&8\\
2&177393&8&1\\
3&177555&8&54\\
4&177771&8&11\\
5&339555&8&6\\
6&339771&8&15\\
7&339933&8&59\\
8&393177&8&0\\
9&393393&8&45\\
10&393555&8&40\\
11&555177&8&52\\
12&555339&8&4\\
13&555393&8&41\\
14&555555&24&84\\
15&555717&8&14\\
16&555771&8&18\\
17&555933&8&24\\
18&717555&8&12\\
19&717717&8&21\\
20&717933&8&19\\
21&771177&8&9\\
22&771339&8&13\\
23&771555&8&16\\
24&933339&8&57\\
25&933555&8&26\\
26&933717&8&17\\
\bottomrule
\end{longtable}

Las sumas de multiplicidades son \(18\cdot18=324\) y
\(24\cdot8+24+108=324\). Para obtener la emisión conjunta se aplica
\eqref{iv:eq:acoplamiento} a cada par de firmas. La inversa por centenas
y decenas demostrada en la sección~\ref{iv:sec:extension} identifica las 468
emisiones y da sus multiplicidades por producto de las dos fibras.
La aplicación posterior \eqref{iv:eq:psi-catalogo} produce las 243
palabras visibles. Así se especifican por separado el dominio inicial,
la emisión, sus fibras y el cociente ternario.

\subsection{Generadores explícitos y cierre de las estrellas}
\label{iv:iv:apendice-estrellas}

Se conserva la matriz \(A_W\) de \eqref{iv:exc:aw} y se forman,
en orden lexicográfico, las palabras
\[
 c(w)=(w,wA_W),\qquad w\in\mathbb F_3^6.
\]
Se agrupan por peso y se retienen los soportes de peso seis.
Esto produce el enumerador \eqref{iv:exc:enumerador} y las 132 hexadas
de \(\mathcal H\). Para cada cuaterna \(B\subset\{1,\ldots,12\}\),
los cuatro pares \(H\setminus B\), con \(B\subset H\in\mathcal H\),
definen \(\sigma_B\): fija \(B\) e intercambia los dos puntos
de cada par. Se obtiene así la familia de 495 estrellas de
la proposición~\ref{iv:exc:m12}.

Los seis generadores usados en su cierre se expresan por sus listas
de imágenes de \(1,\ldots,12\):
\[
\begin{array}{c|rrrrrrrrrrrr}
 &1&2&3&4&5&6&7&8&9&10&11&12\\ \hline
g_1&1&2&3&4&12&11&9&10&7&8&6&5\\
g_2&1&2&3&12&5&10&11&9&8&6&7&4\\
g_3&1&2&3&11&10&6&8&7&12&5&4&9\\
g_4&1&2&12&4&5&9&8&7&6&11&10&3\\
g_5&1&12&3&4&5&8&10&6&11&7&9&2\\
g_6&12&2&3&4&5&7&6&11&10&9&8&1
\end{array}
\]
Cada fila se verifica como una estrella comparándola con las
\(\sigma_B\) ya construidas. Para \(j=1,\ldots,6\), el cierre
se calcula mediante
\[
 G_j^{(0)}=\{\operatorname{id}\},\qquad
 G_j^{(n+1)}
 =G_j^{(n)}\cup
   \{g_i\circ p:1\leq i\leq j,\ p\in G_j^{(n)}\}.
\]
La sucesión creciente se estabiliza en el conjunto finito
\(\operatorname{Sym}(12)\). Como los generadores son involuciones,
su valor estable es exactamente el grupo \(G_j\) que generan.
Con la convención \((g_i\circ p)(a)=g_i(p(a))\), sus órdenes son
\[
\begin{array}{c|rrrrrr}
 j&1&2&3&4&5&6\\ \hline
 |G_j|&2&6&18&360&7920&95040 .
\end{array}
\]
Finalmente se comprueba que cada una de las 495 estrellas pertenece
a \(G_6\) y que todas preservan \(\mathcal H\). La primera inclusión,
junto con la pertenencia de los seis generadores a la familia,
iguala los dos grupos generados. La preservación se comprueba
mediante \(\{\sigma_B(H):H\in\mathcal H\}=\mathcal H\).
Ésta es la enumeración finita utilizada en la prueba de
\ref{iv:exc:m12}; el reconocimiento posterior como \(M_{12}\)
conserva la referencia externa allí declarada.

